\section{Experiments}
\label{sec:results}

\subsection{Setup}

All models are trained and evaluated in IR-SIM, a lightweight 2D
simulator also used by NeuPAN's original benchmarks, ensuring our
LiDAR-based comparison uses an identical sensing modality and
collision model. The robot is modeled as a differential-drive disk
of radius 0.2\,m (approximating TurtleBot3 Burger). Training uses
60--80 epochs of 70 episodes each (300-step cap), Adam optimizer,
batch size 64, on an RTX 5060 GPU. For real-world validation, the
trained policy is deployed without fine-tuning on a TurtleBot3 Burger
equipped with an RPLIDAR C1 and an Orange Pi 5 Pro for on-board
inference.

\subsection{Hard Scenario Benchmark}

To expose failure modes invisible to standard randomized-obstacle
evaluation, we construct four structured scenarios: \textbf{S1
U-trap} (agent inside a concave enclosure, goal behind the opening),
\textbf{S2 Double-U} (two sequential concave traps), \textbf{S3
Narrow-door} (0.45\,m gap, requiring precise alignment), and
\textbf{S5 Symmetric-corridor} (bilaterally symmetric walls inducing
identical left/right LiDAR returns). Each scenario is tested with 4
initial headings $\times$ 3 repetitions (12 trials), using fixed
obstacle layouts to isolate policy capability from environment
randomness.

\subsection{Main Results}

\begin{table}[t]
\centering
\caption{Success rate (\%) across standard generalization
(N=50 trials, fixed seed) and four structured hard scenarios.}
\label{tab:main}
\begin{tabular}{lccccc}
\toprule
Method & Gen. & S1 & S2 & S3 & S5 \\
\midrule
CNNTD3            & 80 & 0   & 100 & 100 & 100 \\
RCPG (GRU)        & 88 & 0   & 0   & 75  & 100 \\
CNNTD3-v2         & 86 & 0   & 100 & 100 & 100 \\
CNNTD3-v3         & 86 & 0   & 75  & 100 & 83  \\
\textbf{CNNTD3-Improved} & 78 & \textbf{100} & 100 & 0 & 100 \\
NeuPAN~\cite{han2025neupan} & 0  & 0   & 0   & 0   & 0   \\
\bottomrule
\end{tabular}
\end{table}

Table~\ref{tab:main} reports results on a 50-trial generalization
test with fixed random seeds and the four hard scenarios. CNNTD3-
Improved is the only method achieving non-zero success on S1 U-trap
(100\%), while all other variants -- including the unmodified
baseline -- fail completely (0\%), confirming that escaping concave
traps requires the targeted curriculum and exploration reward rather
than emerging naturally from standard training. Generalization scores
across variants are statistically overlapping
(78--88\%, individual 95\% confidence intervals span 10--12 points),
indicating that the curriculum modification does not significantly
degrade standard-environment performance despite the architectural
and reward changes.

GRU-based memory (RCPG) shows the inverse pattern: it achieves the
highest generalization score (88\%) and strong narrow-door performance
(75\%) but fails entirely on S2 Double-U (0\%), where path persistence
induced by temporal memory prevents the backtracking required to
escape the second trap. This double-edged effect -- beneficial for
precision tasks, harmful for traps requiring reversal -- is consistent
across all tested scenarios and suggests memory augmentation is not a
general-purpose solution for structured navigation failures.

NeuPAN fails across all five evaluation settings (0\% in every
scenario), including the standard environment using its original
1.6$\times$2.0\,m configuration matched to our scene layout. Diagnostic
analysis (Section~\ref{sec:neupan-analysis}) confirms this is due to
its conservative \texttt{collision\_threshold} mechanism halting motion
when the minimum obstacle distance falls below a fixed bound --
triggered almost immediately given the obstacle density in our
benchmark scenes.

\subsection{Ablation: Exploration Reward Magnitude}

\begin{table}[t]
\centering
\caption{Ablation over exploration reward strength and curriculum
schedule, all evaluated on S1 U-trap.}
\label{tab:ablation}
\begin{tabular}{lcccc}
\toprule
Variant & Explore. bonus & Curr. start & S1 SR & Gen. SR \\
\midrule
CNNTD3 (no curr.) & --   & --       & 0   & 80 \\
CNNTD3-v3          & 0.10 & epoch 20 & 0   & 86 \\
CNNTD3-v2          & 0.15 & epoch 20 & 0   & 86 \\
\textbf{Improved}  & \textbf{0.30} & epoch 10 & \textbf{100} & 78 \\
\bottomrule
\end{tabular}
\end{table}

Table~\ref{tab:ablation} isolates the effect of exploration reward
magnitude. Weaker settings (0.10--0.15) fail to escape U-trap despite
otherwise similar curriculum schedules, while only the strongest
setting (0.30, paired with an earlier curriculum onset at epoch 10)
succeeds. This is not a smooth trade-off curve: intermediate values
produce intermediate generalization cost (78--86\%) but \emph{zero}
intermediate trap-escape benefit, suggesting a threshold effect rather
than a continuous dial between the two objectives. We interpret this
as evidence that escaping a concave trap requires sufficiently strong
anti-stagnation pressure to overcome the locally-optimal "approach and
collide" policy that the base reward otherwise favors near a wall.

\subsection{NeuPAN Failure Analysis}
\label{sec:neupan-analysis}

To rule out robot-size mismatch as the sole cause of NeuPAN's failure,
we retrain its DUNE perception module with TurtleBot3 Burger dimensions
(0.178$\times$0.138\,m, 5000 training epochs) and re-evaluate on a
20\,m corridor scenario with two perpendicular obstacles. With matched
robot geometry, NeuPAN no longer triggers its immediate safety-stop;
instead, it enters a \emph{persistent oscillatory local minimum}
approximately 7.7\,m into the corridor. Across three independent
repeated trials, the robot converges to nearly identical terminal
positions (7.75--7.77\,m, $<$0.02\,m variance), with the last 100
steps showing position standard deviation of 0.013--0.017\,m and
velocity sign reversals in 53--60\% of consecutive steps, while
\texttt{min\_distance} remains well above the 0.05\,m collision
threshold throughout (0.9--1.0\,m). This rules out the safety-stop
mechanism as the cause and instead indicates a genuine failure of the
underlying MPC optimization to converge to a stable forward-moving
solution at this obstacle geometry -- a limitation independent of
parameter tuning, and to our knowledge not previously reported for
NeuPAN in constrained corridor settings.

\subsection{Inference Efficiency}

The CNNTD3 policy (TorchScript, CPU) runs at 0.15\,ms per inference
step on a desktop CPU and 1.65\,ms on the Orange Pi 5 Pro ARM
processor used for on-board deployment -- well within the 100\,ms
budget required for 10\,Hz control. NeuPAN requires 5.3\,ms per step
under comparable conditions, making our policy approximately 3$\times$
faster on edge hardware despite NeuPAN's substantially higher failure
rate in our benchmark.
